% =============================================================
%  ch14_cad_urdf.tex  --  Vorabschritte: vom CAD-Modell zur Simulation
% =============================================================

\chapter{Vorabschritte: vom CAD-Modell zur Simulation}
\label{ch:cad_urdf}

Bevor du simulierst, brauchst du ein \textbf{URDF} (Unified Robot Description
Format) -- die XML-Beschreibung deines Roboters aus \emph{Links} (starre
Körper) und \emph{Joints} (Gelenke dazwischen).

\section{Getrennte STL-Exporte pro Gelenk}
Wenn zwischen Teil\,A und Teil\,B eine Gelenkverbindung besteht, müssen beide
\emph{separat} als Mesh (STL/DAE) exportiert werden -- denn jedes ist ein
eigener Link mit eigenem Koordinatensystem, und das Gelenk dazwischen wird in
der URDF definiert, nicht im CAD-Mesh. Vorgehen:
\begin{enumerate}
\item Im CAD die Gelenkachsen und je Link ein Referenz-Frame festlegen
  (idealerweise Mesh-Ursprung = Gelenkpunkt -- erspart später viel Frame-Frust).
\item \textbf{Jeden Link einzeln} als STL exportieren (Teil\,A, Teil\,B,
  \dots), bevorzugt in seinem eigenen Frame.
\item Masse + Trägheitstensor je Link aus dem CAD übernehmen (für korrekte
  Dynamik).
\item URDF schreiben: Links mit \code{visual}/\code{collision}/\code{inertial};
  Joints mit \code{parent}/\code{child}, \code{origin} (Lage des Gelenks),
  \code{axis} und \code{limit}.
\end{enumerate}

\begin{infobox}[title=Automatisierung -- nicht von Hand schreiben]
Nutze einen CAD-zu-URDF-Exporter statt das XML manuell zu tippen:
\textbf{sw2urdf} (SolidWorks), \textbf{onshape-to-robot} (OnShape) oder
Fusion-360-Plugins. Du wählst Gelenkachsen und Referenz-Frames interaktiv; das
Tool exportiert die Meshes \emph{und} eine konsistente URDF mit korrekten
Frames und Trägheiten.
\end{infobox}

\begin{lstlisting}[style=xml,caption={URDF-Auszug: zwei Links mit Drehgelenk}]
<link name="teilA">
  <visual>
    <geometry>
      <mesh filename="package://my_robot/meshes/teilA.stl"/>
    </geometry>
  </visual>
  <collision>
    <geometry>
      <mesh filename="package://my_robot/meshes/teilA.stl"/>
    </geometry>
  </collision>
  <inertial>
    <mass value="0.45"/>
    <inertia ixx="0.001" iyy="0.001" izz="0.0008"
             ixy="0" ixz="0" iyz="0"/>
  </inertial>
</link>

<joint name="gelenk_AB" type="revolute">
  <parent link="teilA"/>
  <child  link="teilB"/>
  <origin xyz="0 0 0.08" rpy="0 0 0"/>
  <axis   xyz="0 0 1"/>
  <limit lower="-1.57" upper="1.57" effort="12.0" velocity="3.0"/>
</joint>

<link name="teilB"> <!-- visual/collision/inertial wie oben --> </link>
\end{lstlisting}

\section{Gelenktypen}
\begin{itemize}
\item \code{revolute} -- Drehgelenk mit Anschlag (Limits). Der Normalfall für
  Roboterarme.
\item \code{continuous} -- Drehgelenk ohne Anschlag (Räder).
\item \code{prismatic} -- Linearführung.
\item \code{fixed} -- starre Verbindung (z.\,B.\ Kamera an Basis).
\end{itemize}

\section{Getriebeübersetzung in ros2\_control}
Motor + Zykloidgetriebe werden als \emph{ein} angetriebenes Gelenk abstrahiert.
Die Untersetzung steckt in der \textbf{ros2\_control-Transmission}:
\begin{lstlisting}[style=xml,caption={Getriebeübersetzung als SimpleTransmission}]
<transmission name="trans_gelenk_AB">
  <type>transmission_interface/SimpleTransmission</type>
  <joint name="gelenk_AB">
    <hardwareInterface>
      hardware_interface/PositionJointInterface
    </hardwareInterface>
  </joint>
  <actuator name="motor_AB">
    <mechanicalReduction>40</mechanicalReduction>  <!-- 40:1 -->
  </actuator>
</transmission>
\end{lstlisting}
\noindent Die Physik-Engine sieht nur das Ausgangsgelenk. Die Übersetzung
bildet das Verhältnis Motorwinkel $\leftrightarrow$ Gelenkwinkel ab; die
\emph{nach} dem Getriebe verfügbaren Werte trägst du als \code{effort}/
\code{velocity}-Limits am Joint ein (hohes Moment, niedrige Drehzahl). Reales
Restspiel und Reibung kannst du optional über die Joint-Dynamik
(\code{damping}, \code{friction}) annähern.
